\documentclass[11pt]{article}
\usepackage[margin=0.9in]{geometry}
\usepackage{amsmath,amssymb,amsthm,booktabs,mathrsfs}
\usepackage[T1]{fontenc}
\usepackage{lmodern,microtype}
\usepackage[colorlinks=true,linkcolor=blue,urlcolor=blue]{hyperref}
\newtheorem{theorem}{Theorem}[section]
\newtheorem{lemma}[theorem]{Lemma}
\newtheorem{proposition}[theorem]{Proposition}
\newtheorem{corollary}[theorem]{Corollary}
\theoremstyle{remark}
\newtheorem{remark}[theorem]{Remark}
\newcommand{\norm}[1]{\left\|#1\right\|}
\title{Toward a four-dimensional Yang--Mills construction\\
\large Notebook 21: local classical actions through arbitrary refinement}
\author{Research working document}
\date{17 October 2026}
\begin{document}
\maketitle
\begin{abstract}
We localize the actual homogeneous constrained classical action
uniformly through the preceding refinement $M=L^k$. Smooth
physical cells and the normalized fine pairing give uniform
mixed-insertion bounds for the concrete Wilson and corrected
constraint vertices. A contraction in a rooted support algebra
then constructs the minimizer and multiplier as convergent
families of functions of specified background cells. The
action-envelope identity gives strictly supported scalar
coefficients; their Hessians are individually symmetric local
precision pieces. The algebra admits an arbitrary fixed
support-cardinality weight, at the cost of a smaller uniform
background ball. Substitution into the exact fixed-$L$ source
chart gives the localized classical remainder and physical
precision needed by the overlap integration theorem. This
extends the bare small-field integral of Notebook 20 to one
integration of the classical reference action at any earlier
classical depth. Generated quantum interactions and large-field
histories still require an induction theorem.
\end{abstract}

The distinction between the two kinds of depth is explicit:
$S_M$ is obtained by classical minimization, while a full
quantum effective action also contains the corrections from
every earlier fluctuation integral. The theorem here controls
the former and one subsequent small-field integration. All
strict support claims concern the original independent
current-background variables.

